%%%%%%%% ICML 2026 SUBMISSION %%%%%%%%%%%%%%%%%

\documentclass{article}

\usepackage{microtype}
\usepackage{graphicx}
\usepackage{subcaption}
\usepackage{booktabs}

\usepackage{hyperref}

\newcommand{\theHalgorithm}{\arabic{algorithm}}

\usepackage{icml2026}

\usepackage{amsmath}
\usepackage{amssymb}
\usepackage{mathtools}
\usepackage{amsthm}

\usepackage[capitalize,noabbrev]{cleveref}
\crefname{proposition}{Proposition}{Propositions}
\Crefname{proposition}{Proposition}{Propositions}

%%%%%%%%%%%%%%%%%%%%%%%%%%%%%%%%
% THEOREMS
%%%%%%%%%%%%%%%%%%%%%%%%%%%%%%%%
\theoremstyle{plain}
\newtheorem{theorem}{Theorem}[section]
\newtheorem{proposition}[theorem]{Proposition}
\newtheorem{lemma}[theorem]{Lemma}
\newtheorem{corollary}[theorem]{Corollary}
\theoremstyle{definition}
\newtheorem{definition}[theorem]{Definition}
\newtheorem{assumption}[theorem]{Assumption}
\theoremstyle{remark}
\newtheorem{remark}[theorem]{Remark}

\usepackage[textsize=tiny]{todonotes}

\icmltitlerunning{Selective Adversarial Causal Oversight}

\begin{document}

\twocolumn[
  \icmltitle{Selective Adversarial Causal Oversight:\texorpdfstring{\\}{}%
    A Leakage-Free Benchmark and Countermodel-Grounded Verifier under Information Asymmetry}

  \icmlsetsymbol{equal}{*}

  \begin{icmlauthorlist}
    \icmlauthor{Jingbei Niu}{zju}
    \icmlauthor{Zeyuan Tong}{zju}
    \icmlauthor{Yuhang Li}{zju}
  \end{icmlauthorlist}

  \icmlaffiliation{zju}{Zhejiang University, Hangzhou, China}

  \icmlcorrespondingauthor{Yuhang Li}{liyuhangnice666@gmail.com}

  \icmlkeywords{Causal Inference, Large Language Models, Selective Prediction,
    Adversarial Evaluation, Benchmarks, Wise Refusal}

  \vskip 0.3in
]

\printAffiliationsAndNotice{}

\begin{abstract}
LLMs are increasingly used to evaluate causal claims, but real causal oversight
is both \emph{selective} and \emph{adversarial}: claimants may hold private
information, exploit non-identifiability, and present persuasive but incomplete
arguments. We formalize this setting as \emph{selective adversarial causal
oversight}, where a verifier with only public evidence must decide whether a
natural-language claim is \textsc{valid}, \textsc{invalid}, or
\textsc{unidentifiable}. We release a leakage-free benchmark with schema-level
public/gold partitioning, twelve graph families across Pearl's hierarchy,
structurally grounded attacks, a persuasion overlay, and OOD buckets for
mechanism, attack-family, context shift, and paired-flip shifts. We also
propose \emph{countermodel-grounded selective verification}, a four-stage
pipeline that parses the claim, builds an assumption ledger, searches for
observationally compatible countermodels, and returns an abstention-aware
verdict. On the current mainline snapshot, the full system improves the
accuracy/unsafe-acceptance/wise-refusal tradeoff over judge-style baselines,
while a deliberately leaking positive control behaves exactly as our
oracle-leakage proposition predicts. These results illustrate why
public/gold partitioning is a necessary contract for causal oversight
benchmarks.
\end{abstract}

\section{Introduction}
\label{sec:intro}

LLMs are now used as judges, debate moderators, and policy explainers, but
evaluating their \emph{causal} reasoning remains awkward. Existing benchmarks
largely treat causal questions as graded multiple-choice problems
\cite{jin2023cladder,jin2024corr2cause,frohberg2022crass}: the evaluator
inspects an oracle structural causal model, fixes a yes/no answer, and scores
the LLM's prediction against it. This format is convenient but masks two
features that dominate causal oversight in practice.

\textbf{Information asymmetry.} A claimant -- a clinician, policy advocate,
research summarizer, or deceptive agent -- often has access to private
information that the verifier does not: the data-generating process, hidden
confounders, selection mechanisms, or simply domain familiarity. The verifier
sees only public evidence and the claim text. Ignoring this asymmetry
collapses a hard \emph{oversight} problem into an easy
\emph{closed-world recall} problem.

\textbf{Wise refusal.} A core lesson of causal inference is that, given a fixed
public evidence base, many queries are not identifiable
\cite{pearl2009causality,bareinboim2022pearl}. In such cases the correct
behaviour is not to commit to a verdict, but to abstain with a reason. Modern
LLM judges, however, are trained to commit and exhibit sycophancy and
position bias when uncertain
\cite{sharma2024sycophancy,chen2024judge,ye2024justice}. Existing causal
benchmarks rarely measure ``did the model refuse for the right reason'' as a
first-class quantity.

\paragraph{Our task.} We formalize \emph{selective adversarial causal oversight}
(Section~\ref{sec:task}). An attacker generates a natural-language claim with
full knowledge of the gold structural causal model (SCM); a verifier sees only
the public projection of the scenario plus the claim text; and the evaluator
scores the verifier's \emph{three-way} verdict
$\{\textsc{valid}, \textsc{invalid}, \textsc{unidentifiable}\}$ together with a
structured assumption ledger, refusal reason, and witness.

\paragraph{Our benchmark.} We construct a synthetic core of twelve graph
families that explicitly mark family-level identifiability, observationally
equivalent counter-structures, and a finite catalogue of structural attack
modes. A persuasion overlay projects each claim into authority, expert-tone,
confidence, consensus and concealment styles. Splits enforce partitioned
holdouts at the family, mechanism, attack-family, context-shift, and paired-flip
levels, and leakage-prone metadata fields are stripped from the verifier-side
view by construction (Section~\ref{sec:benchmark}). A literature-grounded
subset is reserved as an external sanity check.

\paragraph{Our verifier.} \emph{Countermodel-grounded selective verification}
(Section~\ref{sec:method}) is a four-stage pipeline that (i) parses the claim
into structured slots, (ii) builds an assumption ledger of identifying
preconditions, (iii) searches for observationally compatible counter-SCMs that
disagree with the claim on the target query, and (iv) consults Pearl-style
identification tools as supporting evidence. A fixed decision rule turns these
artifacts into one of the three verdicts plus an abstention reason and a
missing-information specification.

\paragraph{Our claims.} Three propositions
(Section~\ref{sec:theory}) tie selective abstention, countermodel witnesses,
and oracle leakage to formal properties of the task. Empirically
(Section~\ref{sec:experiments}), on the current three-seed,
ten-samples-per-family workspace snapshot we observe that (a) the full system
improves verdict accuracy over \textsc{Direct Judge} from $0.50$ to $0.83$ on
\texttt{test\_iid} and from $0.16$ to $0.85$ on \texttt{test\_ood}, while
driving unsafe acceptance from $0.89/0.97$ to $0.00/0.00$; (b) a deliberate
oracle-leaking positive control reaches $1.00$ accuracy on both dedicated
attack-only splits and therefore remains a useful negative example for any
benchmark that does not enforce a public/gold partition; and (c) removing
\emph{abstention} preserves much of IID raw accuracy ($0.79$) but collapses
wise-refusal recall to $0.00$ on both splits, which is exactly the failure mode
the task is designed to surface. We discuss limitations -- in particular, the
still-modest benchmark scale and the exploratory status of the leakage study --
in Section~\ref{sec:experiments} and Section~\ref{sec:limitations}.

\section{Related Work}
\label{sec:related}

\paragraph{Causal benchmarks for LLMs.}
\textsc{CLadder} \cite{jin2023cladder} introduces a $10$K-sample symbolic
causal benchmark spanning Pearl's three rungs, with ground truth derived from
an oracle inference engine. \textsc{Corr2Cause} \cite{jin2024corr2cause} probes
whether LLMs can convert correlational statements into causal structure; the
authors report near-random performance and a sharp drop under perturbation.
\textsc{CRASS} \cite{frohberg2022crass} targets counterfactual conditionals.
We differ from these benchmarks in three structural ways: (i) we strip oracle
fields from the verifier's view by schema construction rather than by prompt
hygiene, (ii) we elevate \textsc{unidentifiable} to a first-class supervised
label, and (iii) we evaluate not just accuracy but a wise-refusal-aware
metric bundle and a deliberately leaking positive control.

\paragraph{Selective prediction and abstention.}
Classical selective prediction
\cite{geifman2017selective,geifman2019selectivenet} introduces a reject option
to deep classifiers. Recent work re-derives this idea for LLMs and finds that
abstention quality is sensitive to calibration
\cite{guo2017calibration,wen2025abstention}. Our task definition is closest to
this literature in spirit: the verifier explicitly outputs an abstention reason
when public evidence does not determine the query, for example observational
equivalence or missing identifying support.

\paragraph{LLM-as-judge and sycophancy.}
LLM judges show systematic biases: sycophancy
\cite{sharma2024sycophancy}, position and verbosity bias \cite{ye2024justice},
and authority effects \cite{chen2024judge}. Causal claims compound these biases
because the surface plausibility of an explanation is correlated with the
structural correctness of the underlying identification. Our persuasion
overlay is designed to expose exactly that gap.

\paragraph{Scalable oversight and debate.}
\citet{irving2018debate} and \citet{bowman2022measuring} frame oversight as a
game between informed agents and a less-informed judge. Our setting is a
specialization of that frame to causal claims, with a structured rather than
free-form transcript and an explicit abstention option for the judge.

\paragraph{Causal identification and tools.}
We build on canonical identification machinery: do-calculus and the
backdoor/frontdoor criteria \cite{pearl2009causality,shpitser2008complete},
LATE-style instrumental-variable identification \cite{imbens1994late}, and
probability-of-causation results \cite{tian2000probabilities}. Tool execution
in our pipeline uses estimators in the spirit of DoWhy
\cite{sharma2020dowhy}.

\section{Task: Selective Adversarial Causal Oversight}
\label{sec:task}

\subsection{Information Partition}
A scenario is generated as a structural causal model
$\mathcal{M} = (\mathbf{V}, \mathbf{U}, F, P(\mathbf{U}))$ over observed
variables $\mathbf{V}$ and hidden variables $\mathbf{U}$. We expose three
contracted views:

\begin{itemize}
\item \emph{Attacker view}: the gold SCM, the true DAG, hidden variable
  identities, full data, and a generator metadata bundle. The attacker uses
  this to craft a claim that is rhetorically plausible against public
  evidence but structurally constrained.
\item \emph{Verifier view}: a \texttt{PublicCausalInstance} containing only
  observed variables, public observed data, optional named role hints
  (proxy/instrument/mediator/selection), and the natural-language claim.
\item \emph{Evaluator view}: gold + verifier views, used solely to score the
  verdict and audit the witnesses.
\end{itemize}

The contract is enforced at the schema level: the
\texttt{PublicCausalInstance} dataclass strips gold-only fields such as the
true DAG, hidden variables, full SCM, full data, and gold labels, and sanitizes
public text against explicit leakage patterns before the verifier-side view is
emitted.

\subsection{Output and Label Spaces}

For a parsed claim with target $(X, Y)$ over query type
$q \in \{$association, intervention, counterfactual$\}$, the verifier outputs

\begin{itemize}
\item a verdict $\hat{y} \in \{\textsc{valid}, \textsc{invalid},
  \textsc{unidentifiable}\}$,
\item an identification status $s$ that is either \textsc{identified},
  \textsc{contradicted}, or \textsc{underdetermined}, consistent with
  $\hat{y}$,
\item a refusal reason $r$ when $\hat{y} = \textsc{unidentifiable}$,
\item an assumption ledger of typed entries
  (\textsc{supported}, \textsc{contradicted}, \textsc{unresolved}),
\item witness slots for supporting evidence, countermodels, and assumptions,
  and
\item a missing-information spec listing the unsupported assumptions and the
  evidence that would close the identification gap.
\end{itemize}

\subsection{Why the Third Label is Necessary}
\label{sec:third-label}

For three of our twelve graph families (\texttt{l1\_latent\_confounding},
\texttt{l2\_invalid\_iv}, \texttt{l3\_counterfactual\_ambiguity}), there exist
multiple SCMs $\mathcal{M}_1, \mathcal{M}_2$ that are observationally
equivalent but disagree on the target query. Pearl's causal hierarchy theorem
\cite{bareinboim2022pearl} formalizes the general case: layer-$k$ data
underdetermines layer-$(k+1)$ truth almost everywhere. We therefore treat
\textsc{unidentifiable} not as a degraded label but as the unique correct
output when public evidence does not pin down a single query value.

\section{Benchmark}
\label{sec:benchmark}

\subsection{Synthetic Core: Twelve Graph Families}

The benchmark generator instantiates each scenario from a registry of twelve
deterministic graph family blueprints (Table~\ref{tab:families}), distributed
across the three rungs of Pearl's hierarchy (four L1, five L2, three L3
families). Each family is tagged with a family-level
identifiability status (\textsc{identifiable} vs.\
\textsc{potentially\_unidentifiable}), a list of supported gold labels, and a
catalogue of attack modes that match the family's structure.

\begin{table*}[t]
\caption{Twelve graph families across Pearl's hierarchy. Each family fixes
  observable roles (treatment, outcome, context, adjuster, instrument, mediator,
  proxy, selection, latent confounder), a supported label set, and an attack
  catalogue. ``ID'' = structurally identifiable; ``PU'' = potentially
  unidentifiable.}
\label{tab:families}
\begin{center}
\begin{small}
\begin{sc}
\resizebox{\textwidth}{!}{%
\begin{tabular}{lllll}
\toprule
Rung & Family & ID/PU & Supported Labels & Primary Attack Modes \\
\midrule
L1 & latent\_confounding         & PU & invalid, unidentifiable & assoc.\ overclaim, hidden-conf.\ denial \\
L1 & selection\_bias             & PU & invalid, unidentifiable & selection-bias obfuscation \\
L1 & proxy\_disambiguation       & ID & valid, invalid          & assoc.\ overclaim \\
L1 & reverse\_causality          & PU & invalid, unidentifiable & assoc.\ overclaim \\
L2 & valid\_backdoor             & ID & valid, invalid          & invalid-adjustment, heterogeneity overgen.\ \\
L2 & valid\_iv                   & ID & valid, invalid          & heterogeneity overgen.\ \\
L2 & invalid\_iv                 & PU & invalid, unidentifiable & weak-IV-as-valid-IV, exclusion claim \\
L2 & subgroup\_heterogeneity     & PU & invalid, unidentifiable & heterogeneity overgen.\ \\
L2 & frontdoor\_partial\_meas.   & PU & valid, unidentifiable   & unidentifiable disguised as valid \\
L3 & counterfactual\_ambiguity   & PU & unidentifiable          & counterfactual overclaim \\
L3 & mediation\_abduction        & ID & valid, invalid          & function-form manipulation \\
L3 & monotonicity\_cross\_world  & PU & invalid, unidentifiable & cross-world failure, function-form \\
\bottomrule
\end{tabular}
}
\end{sc}
\end{small}
\end{center}
\end{table*}

\subsection{Persuasion Overlay}

A separate overlay layer projects each structurally-fixed claim into a
persuasion-styled variant. The taxonomy mirrors the social-science persuasion
literature \cite{zeng2024johnny} but is deliberately narrowed to five axes:
\textsc{authority}, \textsc{expert\_tone}, \textsc{confidence},
\textsc{consensus}, and \textsc{concealment}. The paper-facing primary axis
uses four of these and reports \textsc{expert\_tone} only in the
taxonomy-complete artifact, to avoid double-counting tone effects with
authority effects in the main table.

\subsection{Witnesses}

Every gold sample carries up to three structured witnesses:
a support witness (why the gold label holds under public evidence), a
countermodel witness (an observationally compatible alternative SCM that
disagrees with the claim on the target query), and an assumption witness
(which identifying assumption is unsupported or contradicted). These witnesses
are first-class evaluation objects rather than narrative decoration: they
ground the decision rule in Section~\ref{sec:method} and the qualitative
witness-faithfulness probe in Section~\ref{sec:experiments}.

\subsection{Splits and Holdouts}

The split manifest enforces non-overlapping \texttt{train}, \texttt{dev},
\texttt{test\_iid}, and \texttt{test\_ood} subsets. The holdout strategy
explicitly tracks several orthogonal axes: graph family, mechanism
(e.g.\ confounding vs.\ collider), attack-family, context-shift domain,
paired-flip pairs, plus lexical/template and variable-renaming controls inside
the generator. This goes beyond the lexical-perturbation OOD studied in prior
work \cite{jin2024corr2cause}, which we view as a useful but insufficient
stress test. Section~\ref{sec:experiments} reports the five paper-facing OOD
buckets that correspond directly to the benchmark contract: graph-family,
mechanism, attack-family, context-shift, and paired-flip OOD.

\subsection{Real-Grounded Subset}

A literature-grounded subset (target size $150$--$300$ cases) is curated from
economics, policy, epidemiology and education observational studies. Each case
records a source citation, a public-evidence summary, the
visible-vs-hidden information contract, the gold label, identifying
assumptions, and a witness note. The synthetic and real-grounded subsets share
the same claim schema, so models trained or tuned on the synthetic core can be
zero-shot evaluated on the real-grounded subset. This subset is part of the
benchmark \emph{design}; the current workspace also includes a $150$-case
human-audit package generator, but not yet dual-annotator agreement results.
Full real-grounded evaluation is therefore reserved for the human-audit phase
(see Section~\ref{sec:limitations}).

\section{Method: Countermodel-Grounded Selective Verification}
\label{sec:method}

The verifier pipeline runs four sequential stages and emits a single
\texttt{VerifierDecision} object.

\paragraph{Stage 1: Claim parsing.}
Given a public scenario and the natural-language claim, the parser extracts
treatment, outcome, query type, polarity, strength, mentioned and implied
assumptions, the rhetorical strategy (e.g.\ an adjustment-sufficiency claim
or an IV appeal), and an abstention-risk flag derived from claim-level cues.

\paragraph{Stage 2: Assumption ledger.}
The ledger initializes one entry per identifying precondition required by the
parsed claim. Status starts as \textsc{unresolved} and is updated by
later stages: tool-backed evidence may flip an entry to \textsc{supported},
contradictory evidence flips it to \textsc{contradicted}. Baseline
assumptions (\texttt{consistency}, \texttt{positivity}) are tracked but not
counted as identifying support.

\paragraph{Stage 3: Countermodel search.}
The countermodel module enumerates structurally-compatible alternative SCMs
for the public scenario and scores each on (a) observational match score, (b)
disagreement with the claim on the target query, (c) which identifying
assumption it triggers. The module returns a chosen candidate, a short
explanation, and a verdict suggestion.

\paragraph{Stage 4: Tool-backed adjudication.}
When no countermodel survives, a tool runner executes the appropriate
identification tools for the rung -- conditional independence and association
checks for L1, backdoor, frontdoor and IV estimators for L2, and SCM-based
counterfactual reasoning plus probability-of-necessity bounds for L3
\cite{tian2000probabilities,sharma2020dowhy}. Tool outputs are stored as a
structured trace that records which ledger entries are supported or
contradicted.

\paragraph{Decision rule.}
The four stages are combined by a fixed rule:
\begin{enumerate}
\item If a strong countermodel survives and the claim's verdict suggestion is
  \textsc{invalid}, return \textsc{invalid}.
\item Else if a countermodel survives and the candidates disagree on the
  target query, return \textsc{unidentifiable} with refusal reason
  \texttt{observational\_equivalence}.
\item Else if any identifying assumption is contradicted by tools, return
  \textsc{invalid}.
\item Else if any core identifying assumption remains unsupported or no
  primary-supporting tool record exists, return \textsc{unidentifiable} with
  a missing-identifying-support reason; when the abstention gate fires on a
  high-risk claim, we return the stricter primary-support variant.
\item Else return \textsc{valid}.
\end{enumerate}
This rule is deterministic, auditable, and is implemented in fewer than $300$
lines of Python; the closed-form schema makes selective verdicts directly
comparable across systems.

\section{Theoretical Guarantees}
\label{sec:theory}

We make three minimal claims that pin down why the task design is necessary
rather than aesthetic.

\begin{proposition}[Abstention Necessity]
\label{prop:abstention}
Let $E$ denote the public evidence and $q$ the target query. Suppose there
exist SCMs $\mathcal{M}_1, \mathcal{M}_2$ that both induce $E$ as their
observational marginal but disagree on $q$. Then any verifier that does not
admit {\normalfont\textsc{unidentifiable}} as an output produces an incorrect
verdict on
at least one of $\mathcal{M}_1, \mathcal{M}_2$.
\end{proposition}
\begin{proof}[Proof sketch]
Both SCMs are public-evidence-equivalent, so the verifier's input distribution
is identical under either generative model. Any committed verdict is therefore
the same under both, but $q(\mathcal{M}_1) \neq q(\mathcal{M}_2)$, so the
verdict is wrong under exactly one.
\end{proof}

\emph{Empirical correspondence.} Section~\ref{sec:experiments}
(\textsc{No-Abstention} ablation) instantiates this proposition: removing the
abstention head retains $0.79$/$0.44$ raw accuracy on IID/OOD but zeroes out
wise-refusal recall.

\begin{proposition}[Countermodel Witness Soundness]
\label{prop:witness}
Suppose the verifier exhibits an SCM $\mathcal{M}'$ that (a) is consistent with
$E$ up to an observational match score $s \ge 1 - \varepsilon$ and (b)
disagrees with the claim on $q$. Then under standard SCM identifiability
\cite{shpitser2008complete}, the claim cannot be uniquely determined by $E$
alone for any $\varepsilon$ such that $\mathcal{M}'$ remains in the equivalence
class.
\end{proposition}
\begin{proof}[Proof sketch]
If $\mathcal{M}'$ matches $E$ to within the equivalence-class tolerance and
disagrees with the claim on $q$, then $E$ is consistent with at least two
distinct values of $q$, contradicting any claim that $E$ alone uniquely
determines $q$. The qualitative shuffle/delete probe in
Section~\ref{sec:experiments} provides exploratory empirical support.
\end{proof}

\begin{proposition}[Oracle Leakage Inflation]
\label{prop:leakage}
Let $V_{\text{public}}$ and $V_{\text{leak}}$ be two verifiers that differ only
in that $V_{\text{leak}}$ has read access to a gold-only field
$g \in \texttt{GOLD\_ONLY\_FIELDS}$. Then
$\mathbb{E}[\text{Acc}(V_{\text{leak}})] \ge
 \mathbb{E}[\text{Acc}(V_{\text{public}})]$, with strict inequality whenever
$g$ is sufficient to recover any portion of the gold label distribution.
\end{proposition}

\emph{Empirical correspondence.} Section~\ref{sec:experiments} (Leakage Study)
instantiates this proposition with a positive-control verifier; the artifact's
\texttt{supports\_warning} field operationalizes the proposition as a
benchmark-level contamination gate.

The role of \cref{prop:leakage} is operational: it justifies our leakage
study (Section~\ref{sec:experiments}) by predicting that any inflation
observed there is, by construction, leakage rather than oversight ability.

\section{Experiments}
\label{sec:experiments}

\subsection{Experimental Setup}
\label{sec:setup}

We report the current workspace snapshot: three seeds with ten samples per
family across all twelve graph families. The main benchmark runs at difficulty
$0.70$ and compares the full \emph{countermodel-grounded} verifier against
judge-style, debate-style, tool-only and refusal-aware baselines. The component
ablation suite replaces one or more pipeline stages with a no-op surrogate. We
also include a dedicated attack-only leakage study at difficulty $0.55$ under
the hidden-information-aware profile. These runs are reproducible and
paper-facing, but they are not a saturation study: the synthetic core is still
modest, the leakage study remains exploratory, and the real-grounded and
human-audit phases are not yet complete. The current workspace also ships a
$150$-case audit package and a pilot cross-model transfer matrix; we treat both
as scope markers rather than headline claims.

\subsection{Main Benchmark}

Table~\ref{tab:main} compares the full \textsc{Countermodel-Grounded}
verifier against compact judge-style and system-style baselines on
\texttt{test\_iid} and \texttt{test\_ood}. We report the three metrics most
central to our task definition: verdict accuracy, unsafe-acceptance rate, and
wise-refusal recall. \textsc{Self-Consistency Judge} matches
\textsc{CoT Judge} exactly in this snapshot and is omitted from the compact
table for readability; the full artifact retains it.

\begin{table*}[t]
\caption{Main benchmark, current mainline snapshot (three seeds, ten
  samples per family, difficulty $0.70$). Mean $\pm$ std across seeds; full
  $95\%$ CIs, calibration metrics, and over-refusal rates are reported in
  the accompanying \texttt{exp\_main\_benchmark.md} artifact.}
\label{tab:main}
\begin{center}
\begin{small}
\begin{sc}
\begin{tabular}{llccc}
\toprule
System & Split & Verdict Acc.\ & Unsafe Accept & Wise Refusal Recall \\
\midrule
\textsc{Direct Judge}          & test\_iid & $0.50 \pm 0.09$ & $0.89 \pm 0.19$ & $0.50 \pm 0.50$ \\
\textsc{Direct Judge}          & test\_ood & $0.16 \pm 0.06$ & $0.97 \pm 0.03$ & $0.18 \pm 0.03$ \\
\textsc{CoT Judge}             & test\_iid & $0.29 \pm 0.08$ & $0.00 \pm 0.00$ & $1.00 \pm 0.00$ \\
\textsc{CoT Judge}             & test\_ood & $0.57 \pm 0.05$ & $0.00 \pm 0.00$ & $0.96 \pm 0.04$ \\
\textsc{Tool Baseline}         & test\_iid & $0.54 \pm 0.12$ & $0.37 \pm 0.26$ & $0.67 \pm 0.58$ \\
\textsc{Tool Baseline}         & test\_ood & $0.36 \pm 0.10$ & $0.42 \pm 0.06$ & $0.58 \pm 0.10$ \\
\textsc{Debate Baseline}       & test\_iid & $0.29 \pm 0.08$ & $0.00 \pm 0.00$ & $1.00 \pm 0.00$ \\
\textsc{Debate Baseline}       & test\_ood & $0.43 \pm 0.04$ & $0.00 \pm 0.00$ & $0.59 \pm 0.07$ \\
\textsc{Refusal-Aware}         & test\_iid & $0.61 \pm 0.10$ & $0.00 \pm 0.00$ & $1.00 \pm 0.00$ \\
\textsc{Refusal-Aware}         & test\_ood & $0.47 \pm 0.09$ & $0.00 \pm 0.00$ & $0.57 \pm 0.06$ \\
\textsc{Countermodel-Grounded} & test\_iid & $0.83 \pm 0.11$ & $0.00 \pm 0.00$ & $1.00 \pm 0.00$ \\
\textsc{Countermodel-Grounded} & test\_ood & $0.85 \pm 0.01$ & $0.00 \pm 0.00$ & $0.79 \pm 0.06$ \\
\bottomrule
\end{tabular}
\end{sc}
\end{small}
\end{center}
\end{table*}

Three observations matter. First, \textsc{Direct Judge} exhibits exactly the
failure mode that motivates the task: extremely high unsafe acceptance
($0.89$ IID, $0.97$ OOD) under public evidence constraints. Second,
\textsc{CoT Judge}, \textsc{Debate Baseline}, and the explicit
\textsc{Refusal-Aware} baseline drive unsafe acceptance to zero, but do so by
over-refusing: on \texttt{test\_ood}, their over-refusal rates are
$0.86$, $0.70$, and $0.59$, respectively, versus $0.07$ for the full system.
Third, \textsc{Countermodel-Grounded} is the best-balanced selective verifier
in this snapshot: it maintains zero unsafe acceptance, preserves strong
wise-refusal recall, and substantially improves verdict accuracy over both
direct and tool-style baselines.

\subsection{Component Ablation: What Drives Wise Refusal?}

Table~\ref{tab:ablation} extends the ablation to four pipeline components,
including a \textsc{No-Abstention} variant that disables the third-label
output entirely. Two patterns are especially important. First,
\textsc{No-Abstention} preserves much of the IID raw accuracy of the full
system ($0.79$ vs.\ $0.83$) but zeroes out wise-refusal recall on both splits,
which is exactly the failure mode the selective task is designed to surface.
Second, removing the countermodel stage hurts both OOD verdict accuracy and OOD
wise-refusal recall ($0.49$ and $0.59$). The current snapshot is not a strict
monotone ranking across every metric: \textsc{No-Tools} slightly exceeds the
full system on OOD raw accuracy ($0.87$ vs.\ $0.85$), but does so with
substantially higher over-refusal ($0.19$ vs.\ $0.07$) and lower
wise-refusal precision ($0.84$ vs.\ $0.93$).

\begin{table}[t]
\caption{Component ablation. Wise-refusal recall is the task-critical
  third-label measure; the full artifact includes precision, over-refusal,
  calibration, and countermodel coverage.}
\label{tab:ablation}
\vskip 0.05in
\begin{center}
\begin{small}
\begin{sc}
\resizebox{\columnwidth}{!}{%
\begin{tabular}{llcc}
\toprule
System & Split & Acc.\ & Wise Refusal Recall \\
\midrule
\textsc{Full}            & iid & $0.83$ & $1.00$ \\
\textsc{Full}            & ood & $0.85$ & $0.79$ \\
\textsc{No-Ledger}       & iid & $0.18$ & $1.00$ \\
\textsc{No-Ledger}       & ood & $0.57$ & $1.00$ \\
\textsc{No-Countermodel} & iid & $0.61$ & $1.00$ \\
\textsc{No-Countermodel} & ood & $0.49$ & $0.59$ \\
\textsc{No-Abstention}   & iid & $0.79$ & $0.00$ \\
\textsc{No-Abstention}   & ood & $0.44$ & $0.00$ \\
\textsc{No-Tools}        & iid & $0.50$ & $1.00$ \\
\textsc{No-Tools}        & ood & $0.87$ & $0.93$ \\
\bottomrule
\end{tabular}
}
\end{sc}
\end{small}
\end{center}
\end{table}

\subsection{Leakage Study}

The current leakage study is explicitly exploratory because it uses only ten
samples per family. We therefore present it as a positive-control sanity check
that the public/gold contract behaves as Proposition~\ref{prop:leakage}
predicts, rather than as a definitive estimate of contamination magnitude.
Full leakage evidence requires increasing samples\_per\_family to at least
$60$.

We instantiate two verifiers on a dedicated attack-only benchmark with the
hidden-information-aware attack profile: a clean public verifier and a
deliberately oracle-leaking control that decodes benchmark-answer supervision
from public metadata. Per \cref{prop:leakage}, any positive gap is
\emph{by construction} leakage inflation. Table~\ref{tab:leakage} shows a gap
of $+0.19$ on \texttt{test\_iid} and $+0.12$ on \texttt{test\_ood}. The
artifact's warning logic is supported on \texttt{test\_ood} but not yet on
\texttt{test\_iid}, so the run only partially supports a formal leakage
warning. Even so, ECE drops from $0.168$ to $0.001$ on \texttt{test\_iid} and
from $0.101$ to $0.001$ on \texttt{test\_ood}, illustrating how leaked
supervision can masquerade as calibration.

\begin{table}[t]
\caption{Leakage study. Clean public verifier vs.\ a positive control that
  reads gold-only fields. Any positive gap is leakage inflation
  (\cref{prop:leakage}). Numbers from the accompanying
  \texttt{exp\_leakage\_study.md} artifact.}
\label{tab:leakage}
\vskip 0.05in
\begin{center}
\begin{small}
\begin{sc}
\begin{tabular}{lccc}
\toprule
Metric (test\_iid) & Clean & Leaking & $\Delta$ \\
\midrule
Accuracy             & $0.81$ & $1.00$ & $+0.19$ \\
Macro-F1             & $0.53$ & $0.67$ & $+0.13$ \\
ECE ($\downarrow$)   & $0.168$ & $0.001$ & $-0.167$ \\
Brier ($\downarrow$) & $0.076$ & $0.00$ & $-0.076$ \\
\midrule
Metric (test\_ood) & Clean & Leaking & $\Delta$ \\
\midrule
Accuracy             & $0.88$ & $1.00$ & $+0.12$ \\
Macro-F1             & $0.56$ & $0.67$ & $+0.10$ \\
ECE ($\downarrow$)   & $0.101$ & $0.001$ & $-0.100$ \\
Brier ($\downarrow$) & $0.049$ & $0.00$ & $-0.049$ \\
\bottomrule
\end{tabular}
\end{sc}
\end{small}
\end{center}
\end{table}

\subsection{Adversarial Robustness}

Table~\ref{tab:adv} reports robustness across four attack-strength regimes:
\textsc{weak}, \textsc{medium}, \textsc{strong}, and a leak-aware
hidden-information profile. On \texttt{test\_ood}, the verifier stays
in a relatively narrow $0.81$--$0.83$ accuracy band across all four regimes
while maintaining zero invalid acceptance. The \texttt{test\_iid} picture is
less stable: the hidden-information-aware regime drops to $0.58 \pm 0.42$,
making it the hardest current IID stress case.

\begin{table}[t]
\caption{Adversarial robustness. Mean $\pm$ std across seeds; full results are
  reported in the accompanying \texttt{exp\_adversarial\_robustness.md}
  artifact.}
\label{tab:adv}
\vskip 0.05in
\begin{center}
\begin{small}
\begin{sc}
\resizebox{\columnwidth}{!}{%
\begin{tabular}{llcc}
\toprule
Strength & Split & Acc.\ & Wise Refusal Recall \\
\midrule
\textsc{Weak}                       & iid & $0.89 \pm 0.19$ & $0.67 \pm 0.58$ \\
\textsc{Weak}                       & ood & $0.81 \pm 0.05$ & $0.65 \pm 0.07$ \\
\textsc{Medium}                     & iid & $0.89 \pm 0.19$ & $0.67 \pm 0.58$ \\
\textsc{Medium}                     & ood & $0.83 \pm 0.04$ & $0.68 \pm 0.09$ \\
\textsc{Strong}                     & iid & $0.89 \pm 0.19$ & $0.67 \pm 0.58$ \\
\textsc{Strong}                     & ood & $0.82 \pm 0.03$ & $0.66 \pm 0.08$ \\
\textsc{Leak-Aware}                 & iid & $0.58 \pm 0.42$ & $0.67 \pm 0.58$ \\
\textsc{Leak-Aware}                 & ood & $0.81 \pm 0.06$ & $0.79 \pm 0.09$ \\
\bottomrule
\end{tabular}
}
\end{sc}
\end{small}
\end{center}
\end{table}

\subsection{OOD Buckets}

The current OOD artifact evaluates the five buckets defined by the benchmark
contract: graph-family, mechanism, attack-family, context-shift, and
paired-flip OOD. The full system remains perfect on graph-family and mechanism
OOD, stays strong on attack-family ($0.77$ accuracy) and context-shift
($0.80$), and drops sharply on paired-flip ($0.42$). All five buckets retain zero unsafe
acceptance in the current snapshot, so paired-flip is best understood as a
structured abstention-and-disambiguation failure mode rather than a simple
unsafe-acceptance failure.

\subsection{Witness Faithfulness (qualitative)}

We run a small witness-faithfulness probe by shuffling or deleting the
countermodel witness in the verifier output and re-scoring. When the
countermodel witness is shuffled across instances, the system's
unidentifiable-awareness collapses on the families that rely on
observational-equivalence reasoning
(\texttt{l1\_latent\_confounding\_family},
\texttt{l3\_counterfactual\_ambiguity\_family}). This indicates that the
witness is load-bearing rather than ornamental. A full quantitative study
(witness shuffle/delete with paired-bootstrap
$95\%$ CIs across all twelve families) is reserved for future work.

\section{Discussion and Limitations}
\label{sec:limitations}

\paragraph{What we are not claiming.}
We are \emph{not} claiming a state-of-the-art result on a saturated benchmark,
nor that LLM judges are uniformly bad at causal reasoning. The claim is
structural: once the task enforces a public/gold contract and rewards correct
abstention, the ranking of systems changes, and a verifier that takes this
shape seriously offers a better accuracy/refusal tradeoff than verifiers that
flatten the problem into a closed-world judgment task.

\paragraph{Limitations.}
The synthetic core is, by construction, a model of the world rather than the
world itself. The real-grounded subset is small, and full human dual-annotation
is in progress. The current mainline snapshot uses ten samples per family, so
bucket-level estimates can still have wide CIs; the dedicated leakage study is
explicitly exploratory; and the hidden-information-aware IID robustness slice is
not yet stable enough to support strong quantitative claims. Larger runs and
completed human audit will be needed before any broad statement about absolute
performance is defensible. The existing transfer matrix is likewise a pilot:
it shows that family-swapped attacker/verifier profiles are testable inside the
benchmark contract, but it is still too small to support model-family
conclusions. We document these caveats in the artifact and in this paper.

\paragraph{Future work.}
Two priorities: (i) closing the real-grounded subset with dual annotation
and Cohen's $\kappa$ on the four primary audit fields; (ii) scaling the
current transfer pilot into a full-budget cross-model study across verifier
families and sizes, which will let us connect the structural conclusions of
this paper to specific model regimes.

\section{Conclusion}

Selective adversarial causal oversight is the right unit of analysis for
LLM-based causal evaluation. By modelling information asymmetry, treating
\textsc{unidentifiable} as a first-class label, partitioning public from gold
at the schema level, and grounding the verifier in countermodel witnesses,
we obtain a benchmark and a method whose comparison shape is meaningful even
under current mainline budgets. The next step is scale: more samples, more graph
families, real-grounded human-audited cases, and cross-model transfer.

\section*{Software and Data}
The benchmark generator, verifier pipeline, and experiment runners are
released alongside the paper as a reproducibility package. Each experiment
emits a configuration JSON, a seed list, raw per-sample predictions, an
aggregated metrics file, $95\%$ CIs, paired-significance tests, and a
markdown summary. The leakage study additionally emits a leakage-warning
field that downstream consumers can use as a contamination check.

\section*{Impact Statement}
This paper presents work whose goal is to advance the field of Machine
Learning. There are many potential societal consequences of our work, none
which we feel must be specifically highlighted here. We note in particular
that our adversarial framing is intended for evaluation and oversight rather
than deployment: the persuasion overlay is a stress-test artifact for
verifiers, not a recipe for deceptive content generation, and the benchmark
release does not include any leakage-decoder code beyond what is required to
reproduce the negative-control leakage study.

\nocite{pearl2009causality,pearl2018book,bareinboim2022pearl}
\nocite{jin2023cladder,jin2024corr2cause,frohberg2022crass}
\nocite{geifman2017selective,geifman2019selectivenet}
\nocite{sharma2024sycophancy,chen2024judge,ye2024justice}
\nocite{irving2018debate,bowman2022measuring}
\nocite{guo2017calibration,wen2025abstention}
\nocite{tian2000probabilities,sharma2020dowhy,imbens1994late,shpitser2008complete}
\nocite{zeng2024johnny,lanham2023measuring}

\bibliography{example_paper}
\bibliographystyle{icml2026}

%%%%%%%%%%%%%%%%%%%%%%%%%%%%%%%%%%%%%%%%%%%%%%%%%%%%%%%%%%%%%%%%%%%%%%%%%%%%%%%
% APPENDIX
%%%%%%%%%%%%%%%%%%%%%%%%%%%%%%%%%%%%%%%%%%%%%%%%%%%%%%%%%%%%%%%%%%%%%%%%%%%%%%%
\newpage
\appendix
\onecolumn
\section{Schema Excerpts}
\label{app:schema}

The following is an abridged view of the leakage-control contract. The full
\texttt{benchmark/schema.py} file contains additional sanitizers and
serialization helpers; only the contract-level constants are reproduced here.

\begin{verbatim}
GOLD_ONLY_FIELDS = (
    "true_dag", "hidden_variables", "true_scm",
    "full_data", "ground_truth", "gold_label",
)

VERIFIER_VISIBLE_FIELDS = (
    "scenario_id", "description", "variables", "proxy_variables",
    "selection_mechanism", "observed_data", "data",
    "causal_level", "metadata",
)

class VerdictLabel(str, Enum):
    VALID = "valid"
    INVALID = "invalid"
    UNIDENTIFIABLE = "unidentifiable"

class IdentificationStatus(str, Enum):
    IDENTIFIED = "identified"
    CONTRADICTED = "contradicted"
    UNDERDETERMINED = "underdetermined"
\end{verbatim}

\section{Decision Rule (Pseudocode)}
\label{app:decision}

\begin{verbatim}
def decide_verdict(parsed, ledger, cm, tool_trace):
    # Stage 1: strong invalid countermodel
    if cm.found and cm.suggests_invalid:
        return INVALID

    # Stage 2: observational equivalence + query disagreement
    if cm.found and cm.query_disagreement:
        return UNIDENTIFIABLE(reason="observational_equivalence")

    # Stage 3a: contradicted assumption
    if any(e.contradicted for e in ledger):
        return INVALID

    # Stage 3b: missing or unsupported identification
    if any(e.unsupported for e in ledger) \
       or not tools_support_claim(tool_trace) \
       or (needs_explicit_support and no_supported_id):
        return UNIDENTIFIABLE(reason="missing_identifying_support")

    # Stage 3c: abstention gate on high-risk valid claims
    if abstention_gate_reasons(parsed, ledger, tool_trace):
        return UNIDENTIFIABLE(reason="missing_primary_identifying_support")

    # Stage 4: support
    return VALID
\end{verbatim}

\section{Persuasion Overlay Examples}
\label{app:persuasion}

A claim with structural label \textsc{unidentifiable} can be projected into
five persuasion styles. The taxonomy-complete artifact includes
\textsc{expert\_tone}; the paper-facing primary axis collapses to the four
remaining styles to avoid double-counting tone effects.

\begin{small}
\begin{itemize}
\item \textsc{Authority}: ``As established by the foundational paper of the
  field, the effect of $X$ on $Y$ is positive.''
\item \textsc{Confidence}: ``It is unambiguous, on any reading of the data,
  that $X$ causes $Y$.''
\item \textsc{Consensus}: ``Every recent meta-analysis agrees that $X$
  raises $Y$.''
\item \textsc{Concealment}: omits any mention of the latent confounder or
  selection mechanism that the public evidence cannot resolve.
\item \textsc{Expert tone}: heavy technical jargon without a corresponding
  identification argument.
\end{itemize}
\end{small}

\section{Reproducibility Notes}
\label{app:repro}

Each experiment runner emits a fixed bundle of artifacts:

\begin{verbatim}
<run>_config.json
<run>_seed_list.json
<run>_raw_predictions.jsonl
<run>_aggregated_metrics.json
<run>_ci.json
<run>.md   (paper-facing markdown summary)
\end{verbatim}

The aggregator reports seed-mean estimates and paired bootstrap $95\%$ CIs;
the significance file applies Holm-Bonferroni correction within and across
splits. The leakage runner additionally writes a
\texttt{supports\_warning} field that downstream consumers can use as a
hard gate on contamination, in line with \cref{prop:leakage}.

%%%%%%%%%%%%%%%%%%%%%%%%%%%%%%%%%%%%%%%%%%%%%%%%%%%%%%%%%%%%%%%%%%%%%%%%%%%%%%%
\end{document}
